% ============================================================
%  USAGE: paste this file's content into your group report
%  between other sections, OR compile standalone with:
%      xelatex report_integration_and_levels.tex
%
%  If compiling standalone, keep the preamble below.
%  If pasting into a group report, remove everything above
%  and including \begin{document}, and remove \end{document}.
% ============================================================
\documentclass[12pt,a4paper]{article}

\usepackage{fontspec}
\setmainfont{Times New Roman}
\setsansfont{Arial}
\setmonofont[Scale=0.88]{Courier New}
\usepackage[margin=2.5cm]{geometry}
\usepackage{amsmath,amssymb}
\usepackage{booktabs}
\usepackage{enumitem}
\usepackage{listings}
\usepackage{xcolor}
\usepackage{hyperref}
\usepackage{float}
\usepackage{fancyhdr}
\setlength{\headheight}{15pt}

\definecolor{codegray}{rgb}{0.5,0.5,0.5}
\definecolor{codeblue}{rgb}{0.0,0.3,0.7}
\definecolor{codeorange}{rgb}{0.7,0.4,0.0}
\definecolor{codebg}{rgb}{0.97,0.97,0.97}

\lstdefinestyle{pycode}{
  language=Python,
  basicstyle=\ttfamily\small,
  keywordstyle=\color{codeblue}\bfseries,
  commentstyle=\color{codegray}\itshape,
  stringstyle=\color{codeorange},
  backgroundcolor=\color{codebg},
  numbers=left, numberstyle=\tiny\color{codegray},
  breaklines=true, frame=single, framesep=4pt,
  xleftmargin=1.8em, tabsize=4, showstringspaces=false,
  captionpos=b,
}
\lstset{style=pycode}

\pagestyle{fancy}
\fancyhf{}
\lhead{\small XiangQi Chess HCMUT --- Project Report}
\rhead{\small Model Integration and Agent Difficulty Levels}
\cfoot{\thepage}
\renewcommand{\headrulewidth}{0.4pt}
\hypersetup{colorlinks=true, linkcolor=blue!60!black}

\begin{document}

% ============================================================
% NOTE TO GROUP: replace the section number below to match
% your report structure, e.g. \section{...} or \subsection{...}
% ============================================================

\section{Policy Model Integration and Agent Difficulty Levels}
\label{sec:integration}

This section describes the design and implementation of two closely
related contributions: (1) integrating the policy-based ML model into
the existing \texttt{GameLoop} so that all game modes share the same
backend pipeline, and (2) defining a three-level difficulty system for
the ML agent and exposing it through the user interface.

% ============================================================
\subsection{Background and Design Goals}
\label{sec:integration-background}

Before this milestone, the project had a functional rule engine,
a headless \texttt{GameLoop}, and an \texttt{MLAgent} stub that used a
dummy heuristic model.  The following items were not yet implemented:

\begin{itemize}
  \item Named difficulty levels for the ML agent.
  \item Game modes that involve the ML agent in the UI (Human vs ML,
        ML vs Random, ML vs Search).
  \item A level-selection panel for the ML agent in the Pygame menu.
  \item A headless evaluation harness for ML vs Minimax / AlphaBeta.
\end{itemize}

The guiding design principle throughout is \textbf{single source of
truth}: the same \texttt{GameLoop}, the same \texttt{legal\_moves}
function, and the same \texttt{assert\_legal\_move} guard are reused
for every mode — headless or UI, search agent or ML agent.

% ============================================================
\subsection{GameLoop: Headless Architecture}
\label{sec:gameloop}

\texttt{game/game\_loop.py} is the single backend turn engine with no
Pygame dependency.  Its public interface is:

\begin{lstlisting}[caption={GameLoop public API}]
# One complete hidden match
result = run_headless_game(red_agent, black_agent, max_turns=200)

# Step-by-step (used by the Pygame main loop)
loop = GameLoop(red_agent, black_agent, max_turns=200)
result = loop.step()   # returns None while match continues
\end{lstlisting}

For each turn \texttt{step()} performs the following operations in order:

\begin{enumerate}
  \item Check terminal state (\texttt{result\_if\_terminal}).
  \item Check \texttt{ply\_count >= max\_turns}.
  \item Select the current agent by \texttt{state.side\_to\_move}.
  \item Call \texttt{agent.select\_move(state.clone())} --- a clone is
        passed so model inference cannot mutate the authoritative board.
  \item Validate the returned move with \texttt{assert\_legal\_move}.
  \item Apply the move and update the threefold-repetition counter.
  \item Return \texttt{None} (continue) or a \texttt{GameLoopResult}
        with \texttt{winner}, \texttt{reason}, and full \texttt{history}.
\end{enumerate}

The Pygame main loop in \texttt{main.py} reuses the same
\texttt{assert\_legal\_move} call and \texttt{state.clone()} pattern
for consistency with the headless path.

% ============================================================
\subsection{ML Agent: Inference Pipeline}
\label{sec:inference}

\texttt{MLAgent} implements the same \texttt{BaseAgent} interface
(\texttt{select\_move(state) -> Optional[Move]}) as every other agent.
On each call it executes five steps:

\begin{enumerate}
  \item \textbf{Filter legal moves.}
        \texttt{legal\_moves(state)} is called and each move is
        re-validated by \texttt{is\_legal\_move} as a defensive check.
  \item \textbf{Encode the board.}
        \texttt{state.to\_tensor(canonical=True)} produces a
        $(15, 10, 9)$ array.  Channels~0--6 are the seven piece types
        for the current player; channels~7--13 for the opponent;
        channel~14 is the side-to-move flag.  When the current player
        is BLACK the board is rotated $180^\circ$ and colors are swapped
        so the model always sees the board from the same perspective.
  \item \textbf{Forward pass.}
        \texttt{model.forward(state\_tensor)} returns an 8\,100-length
        policy vector.  The policy index for a move is:
        \begin{equation}
          \text{index}(src, dst)
          = \bigl(r_{src} \times 9 + c_{src}\bigr) \times 90
          + \bigl(r_{dst} \times 9 + c_{dst}\bigr).
          \label{eq:policy}
        \end{equation}
  \item \textbf{Map scores to legal moves.}
        \texttt{canonical\_move\_to\_policy\_index(mv, side\_to\_move)}
        applies the coordinate flip for BLACK moves before reading each
        legal move's score from the policy vector.
  \item \textbf{Select the best move.}
        The legal move with the highest score is returned (argmax), with
        an optional epsilon-greedy override for the Easy level
        (Section~\ref{sec:levels}).
\end{enumerate}

\subsubsection{Model Loading}

\texttt{MLAgent.\_load\_model} resolves the model source at construction
time (Table~\ref{tab:model-sources}).

\begin{table}[H]
\centering
\caption{Supported model sources for \texttt{MLAgent}.}
\label{tab:model-sources}
\begin{tabular}{lll}
\toprule
\textbf{\texttt{model\_path}} & \textbf{Loader} & \textbf{Notes} \\
\midrule
\texttt{None} & \texttt{DummyPolicyModel()} &
  Heuristic placeholder; 8\,100-dim output \\
\texttt{*.json} & \texttt{DummyPolicyModel(**cfg)} &
  JSON-configurable dummy weights \\
\texttt{*.pt} / \texttt{*.pth} & \texttt{TorchPolicyAdapter} &
  Trained PyTorch checkpoint \\
\bottomrule
\end{tabular}
\end{table}

\texttt{TorchPolicyAdapter} wraps inference in \texttt{torch.no\_grad()},
adds the batch dimension, and returns a plain Python list so the rest of
\texttt{MLAgent} is framework-agnostic.

\subsubsection{DummyPolicyModel}

Before a trained checkpoint is available, \texttt{DummyPolicyModel}
exercises the full inference path by computing a heuristic 8\,100-dim
policy vector from the encoded tensor.  For each source--destination
pair the score is:

\begin{equation}
  s(src, dst) =
  V_{\text{piece}}(src)
  + B_{\text{cap}} \cdot \mathbf{1}[\text{enemy at }dst]
  - \lambda_c \cdot d_{\text{center}}(dst)
  + \lambda_f \cdot \Delta r_{\text{fwd}}
  \label{eq:dummy}
\end{equation}

with $B_{\text{cap}} = 100$, $\lambda_c = 1.0$, $\lambda_f = 0.25$.
Pairs that have no own piece at the source, or that land on a friendly
piece, receive score $-10^6$ and are never selected.

The critical property is that the output shape (8\,100) and coordinate
convention (canonical) are identical to a trained model's output, so
replacing \texttt{DummyPolicyModel} with a trained \texttt{.pth}
checkpoint requires only passing a different \texttt{model\_path} to
\texttt{MLAgent} --- no other code changes.

\subsubsection{Repetition Heuristics}

When the \textit{Hard} level is active, scores of moves that could
cause threefold repetition or simple backtracking are penalized after
the forward pass:

\begin{lstlisting}[caption={Repetition penalty applied after forward pass}]
for i, mv in enumerate(moves):
    if is_simple_backtrack(state, mv):
        scores[i] -= BACKTRACK_PENALTY        # 1500 pts

    undo = state.apply_move(mv)
    try:
        key = game_loop_position_key(state)
        if visit_counts.get(key, 0) + 1 >= 3:
            scores[i] -= THREEFOLD_PENALTY    # 500 000 pts
    finally:
        state.undo_move(undo)
\end{lstlisting}

\texttt{is\_simple\_backtrack} fires when the agent's own previous move
is exactly reversed in the current candidate.  The threefold penalty is
large enough to override all policy scores.

% ============================================================
\subsection{Difficulty-Level System}
\label{sec:levels}

\subsubsection{Motivation}

A single set of parameters does not produce a useful spread of
playing strength.  Three difficulty levels are defined to: (1)~let
beginners play against a weak opponent, (2)~provide a fair baseline
for comparison with the search-based agents, and (3)~give a strongest
configuration for benchmarking.

\subsubsection{Level Definitions}

The levels are declared in the module-level dictionary
\texttt{ML\_AGENT\_LEVELS}:

\begin{lstlisting}[caption={ML agent level configuration}]
ML_AGENT_LEVELS = {
    "Easy":   {"epsilon": 0.4,  "apply_repetition_heuristics": False},
    "Medium": {"epsilon": 0.0,  "apply_repetition_heuristics": False},
    "Hard":   {"epsilon": 0.0,  "apply_repetition_heuristics": True},
}
\end{lstlisting}

Table~\ref{tab:levels} summarises the effect of each level.

\begin{table}[H]
\centering
\caption{ML agent difficulty levels and their behavioral parameters.}
\label{tab:levels}
\begin{tabular}{lccl}
\toprule
\textbf{Level} & \textbf{$\varepsilon$} &
\textbf{Rep.\ heuristics} & \textbf{Behavioral description} \\
\midrule
Easy   & 0.40 & Off &
  40\% of moves are chosen uniformly at random; \\
  & & & 60\% follow the policy argmax.
  Weak and unpredictable. \\[3pt]
Medium & 0.00 & Off &
  Pure policy argmax.  Can loop when \\
  & & & no capture is available (no anti-repetition). \\[3pt]
Hard   & 0.00 & On  &
  Policy argmax with heavy penalties \\
  & & & for repetition and backtracking. \\
\bottomrule
\end{tabular}
\end{table}

\subsubsection{Epsilon-Greedy Selection (Easy Level)}

With probability $\varepsilon$ the agent samples a legal move uniformly
at random rather than following the policy.  The random branch is
seeded via a \texttt{random.Random} instance for reproducibility.

\begin{lstlisting}[caption={Epsilon-greedy branch in \texttt{select\_move}}]
if self.epsilon > 0.0 and self.rng.random() < self.epsilon:
    return self.rng.choice([mv for mv, _ in move_scores])

# Normal path: argmax
best_move, _ = max(move_scores, key=lambda item: item[1])
return best_move
\end{lstlisting}

\subsubsection{Agent Construction}

The level is specified at construction time:

\begin{lstlisting}[caption={Constructing \texttt{MLAgent} at a chosen level}]
from agents.ml_agent import MLAgent
from core.rules import Color

agent = MLAgent(player_id=Color.RED, level="Hard")
# Resolved attributes:
# agent.name    -> "MLAgent(Hard)"
# agent.epsilon -> 0.0
# agent.apply_repetition_heuristics -> True
\end{lstlisting}

Individual parameters (\texttt{epsilon},
\texttt{apply\_repetition\_heuristics}) can still be overridden
explicitly for ablation experiments, independent of the named level.

% ============================================================
\subsection{New Game Modes}
\label{sec:modes}

Three new game modes were added to the menu and backend
(Table~\ref{tab:modes}), alongside the three original modes.

\begin{table}[H]
\centering
\caption{All six supported game modes after this milestone.}
\label{tab:modes}
\begin{tabular}{lll}
\toprule
\textbf{Mode} & \textbf{Red side} & \textbf{Black side} \\
\midrule
Human vs AI    & Human input        & SearchAgent (AI level) \\
AI vs AI       & SearchAgent (Red)  & SearchAgent (Black) \\
AI vs Random   & SearchAgent        & RandomAgent \\
Human vs ML    & Human input        & MLAgent (ML level) \textbf{[new]} \\
ML vs Random   & MLAgent (ML level) & RandomAgent        \textbf{[new]} \\
ML vs Search   & MLAgent (ML level) & SearchAgent (AI level) \textbf{[new]} \\
\bottomrule
\end{tabular}
\end{table}

In \texttt{main.py}, all modes are handled in a single dispatch function:

\begin{lstlisting}[caption={\texttt{\_build\_agents} dispatch (new branches)}]
def _build_agents(mode, level, red_level=None,
                  black_level=None, ml_level="Hard"):
    if mode == "Human vs ML":
        return Color.RED, None,
               _build_ml_agent(Color.BLACK, ml_level)
    if mode == "ML vs Random":
        return None,
               _build_ml_agent(Color.RED, ml_level),
               RandomAgent(player_id=Color.BLACK)
    if mode == "ML vs Search":
        return None,
               _build_ml_agent(Color.RED, ml_level),
               _build_search_agent(level, Color.BLACK)
    # ... original modes unchanged
\end{lstlisting}

% ============================================================
\subsection{UI Level Selector}
\label{sec:ui-level}

The Pygame menu (\texttt{ui/menu.py}) was extended with an
\textbf{ML Level} panel that appears dynamically depending on the
selected mode (Table~\ref{tab:panels}).

\begin{table}[H]
\centering
\caption{Level panel shown per game mode.}
\label{tab:panels}
\begin{tabular}{ll}
\toprule
\textbf{Mode} & \textbf{Level panel displayed} \\
\midrule
Human vs AI, AI vs Random & Single AI Level selector \\
Human vs ML, ML vs Random & Single ML Level selector \\
ML vs Search & Two side-by-side: ML Level + AI Level \\
AI vs AI & Two columns: Red AI level + Black AI level \\
\bottomrule
\end{tabular}
\end{table}

Two helper methods on \texttt{Menu} determine which panels to render:

\begin{lstlisting}[caption={Panel selection logic in \texttt{Menu}}]
def _mode_uses_ai_level(self):
    return self.selected_mode in (
        "Human vs AI", "AI vs Random", "ML vs Search")

def _mode_uses_ml_level(self):
    return self.selected_mode in (
        "Human vs ML", "ML vs Random", "ML vs Search")
\end{lstlisting}

The selected ML level is stored in \texttt{menu.selected\_ml\_level}
and passed to \texttt{\_build\_agents} when the user starts the game:

\begin{lstlisting}[caption={Passing ML level from menu to agent construction}]
human_color, red_agent, black_agent = _build_agents(
    menu.selected_mode,
    menu.selected_level,          # AI (search) level
    red_level   = menu.selected_red_level,
    black_level = menu.selected_black_level,
    ml_level    = menu.selected_ml_level,   # <-- new
)
\end{lstlisting}

% ============================================================
\subsection{Headless Evaluation: ML vs Minimax / AlphaBeta}
\label{sec:headless-eval}

\texttt{evaluation/headless\_match.py} was extended with a second
simulation function for batch evaluation of ML against search agents:

\begin{lstlisting}[caption={Headless evaluation API}]
records = run_ml_vs_search(
    games=10, max_turns=160,
    ml_level="Hard",
    search_level="Easy",          # depth preset for SearchAgent
    search_algorithm="minimax",   # or "alphabeta"
)
\end{lstlisting}

Each record is a \texttt{HeadlessMatchRecord} containing agent names,
winner, reason (\textit{checkmate}, \textit{stalemate}, etc.),
ply count, and elapsed time.  Results can be exported to JSON:

\begin{verbatim}
python evaluation/headless_match.py \
  --mode ml-vs-search \
  --games 10 --ml-level Hard \
  --search-level Easy --search-algorithm minimax \
  --json-out evaluation/results/ml_vs_minimax.json
\end{verbatim}

Table~\ref{tab:results} shows representative results from smoke tests.

\begin{table}[H]
\centering
\caption{Headless smoke-test results (1 game each, \texttt{max\_turns=60}, seed~42).}
\label{tab:results}
\begin{tabular}{llll}
\toprule
\textbf{Red (MLAgent level)} & \textbf{Black} &
\textbf{Outcome} & \textbf{Reason} \\
\midrule
Easy   & RandomAgent       & Draw       & max\_turns\_reached \\
Medium & RandomAgent       & Draw       & stalemate \\
Hard   & RandomAgent       & Draw       & stalemate \\
Hard   & Easy Minimax (d=1) & Black wins & checkmate \\
\bottomrule
\end{tabular}
\end{table}

The dummy policy model losing to depth-1 Minimax is expected: the
heuristic model is not calibrated for endgame sequences.  Replacing
\texttt{DummyPolicyModel} with a trained checkpoint requires no changes
to the evaluation pipeline.

% ============================================================
\subsection{Safety and Correctness Guarantees}
\label{sec:safety}

Regardless of model quality, the following invariants are maintained:

\begin{enumerate}[leftmargin=*]
  \item \textbf{MLAgent only scores legal moves.}
        The model cannot cause the agent to select an illegal move
        because its output is only read for moves already validated
        by \texttt{legal\_moves(state)}.
  \item \textbf{GameLoop validates every move before applying it.}
        \texttt{assert\_legal\_move} raises \texttt{ValueError} if an
        agent returns an illegal move, so a buggy model cannot corrupt
        game state.
  \item \textbf{State isolation.}
        Agents receive \texttt{state.clone()}.
        Model inference cannot mutate the authoritative board.
  \item \textbf{Graceful fallback on bad model output.}
        Wrong output shapes, NaN/Inf values, and forward-pass exceptions
        all fall back to neutral scores (0.0 per legal move) rather than
        crashing the game.
\end{enumerate}

\end{document}
